%% ============================================================
%%  backend-ai-es/experience-es.tex  —  Experiencia (Backend + Node.js + Java + IA)
%%
%%  Enfoque: microservicios Java, APIs REST con Node.js/NestJS,
%%  orquestación de agentes IA, arquitecturas event-driven, Kafka.
%% ============================================================

\section{Experiencia Profesional}

%% ── 1. IDM Technology / Scotiabank  (Marzo 2026 – Presente) ──
\cvEntryClient%
  {Desarrollador Backend Senior}%
  {IDM Technology}%
  {Scotiabank Peru}%
  {Plataforma de Pagos QR Dinámico}%
  {Marzo 2026 -- Presente \enspace\textbar\enspace Contrato 6 meses}

\begin{cvItems}
  \cvItem{Desarrollo de microservicios backend reactivos cloud-native (\textbf{Java 25, Spring Boot 4,
    Spring WebFlux}) para generación de QR dinámico, procesamiento de pagos y persistencia de
    transacciones mediante \textbf{REST APIs}.}
  \cvItem{Diseño de \textbf{servicios backend de integración} entre \textbf{Google Cloud Platform},
    \textbf{Apigee} y sistemas mainframe IBM AS/400 a través de DataQueues.}
  \cvItem{Autenticación bancaria (\textbf{Spring Security}) y optimización de capas de persistencia
    \textbf{SQL Server} para cargas transaccionales de alto rendimiento.}
  \cvItem{Desarrollo de backends de automatización interna con \textbf{Node.js} y \textbf{Python}
    para optimizar flujos operativos del banco.}
  \cvItem{Pipelines \textbf{Jenkins CI/CD} y calidad con \textbf{SonarCloud}; resolución de
    vulnerabilidades de penetration testing en equipos Agile Scrum.}
  \cvItem{\textbf{Desarrollo asistido por IA} (Claude Code, GitHub Copilot) para acelerar
    scaffolding backend y revisión de código en todos los workstreams.}
\end{cvItems}

{\footnotesize\color{cvMuted}\textit{Stack:}~%
Java 25, Spring Boot 4, Spring WebFlux, Programación Reactiva, REST APIs, SQL Server,
IBM AS/400, DataQueues, Spring Security, GCP, Apigee, Node.js, Python, Jenkins, SonarCloud,
Claude Code, GitHub Copilot.}

\vspace{2pt}

%% ── 2. DaCodes — AgentCloud  (Enero 2025 – Febrero 2026) ────
\cvEntryClient%
  {Ingeniero Backend Líder — Plataforma de Agentes IA}%
  {DaCodes}%
  {HealthCare}%
  {AgentCloud}%
  {Enero 2025 -- Febrero 2026}

\begin{cvItems}
  \cvItem{Arquitectura y desarrollo backend de \textbf{AgentCloud}, plataforma empresarial de
    \textbf{orquestación de agentes IA} para clientes Healthcare en Estados Unidos.}
  \cvItem{Diseño del \textbf{motor de flujos agénticos}: coordinación event-driven de agentes IA
    con \textbf{Apache Kafka} para ejecución de workflows multi-paso en tiempo real.}
  \cvItem{Microservicios cloud-native en \textbf{AWS} (Lambda, EKS, SQS) con \textbf{Java /
    Spring Boot}; contenedorizados con Docker y orquestados con Kubernetes.}
  \cvItem{Implementación de \textbf{REST APIs} para configuración dinámica de agentes, integración
    de LLMs y ejecución de pipelines agénticos; \textbf{ingeniería de prompts} en servicios backend.}
  \cvItem{Uso de \textbf{agentes IA} (Claude Code, GitHub Copilot, Cursor) para scaffolding backend,
    revisiones automatizadas e ingeniería de prompts dentro de la plataforma.}
  \cvItem{Liderazgo de equipo distribuido en ceremonias Agile; \textbf{CI/CD} con GitHub Actions y
    estándares de calidad de código.}
\end{cvItems}

\vspace{2pt}

%% ── 3. ExSquared / SambaSafety — Qorta  (Marzo 2023 – Dic 2024)
\cvEntryClient%
  {Ingeniero Backend Senior}%
  {ExSquared}%
  {SambaSafety}%
  {Qorta}%
  {Marzo 2023 -- Diciembre 2024}

\begin{cvItems}
  \cvItem{Construcción de \textbf{REST APIs} escalables con \textbf{NestJS / Node.js} para Qorta,
    plataforma de monitoreo de riesgo de conductores de SambaSafety en Norteamérica.}
  \cvItem{Integración de backends \textbf{Node.js} con proveedores de registros vehiculares (MVR)
    en más de \textbf{50 jurisdicciones de EE.UU.} para scoring de riesgo en tiempo real.}
  \cvItem{Implementación de flujos de procesamiento de datos con \textbf{Kafka} y \textbf{Redis};
    optimización de queries \textbf{PostgreSQL} para grandes volúmenes de datos de conductores.}
  \cvItem{Mantenimiento de cobertura de pruebas superior al \textbf{85\%} con Jest; colaboración
    en equipos Agile Scrum distribuidos con ingenieros en EE.UU.}
\end{cvItems}

\vspace{2pt}

%% ── 4. Globant / Banco GNB Paraguay  (Junio 2022 – Mayo 2023) ──
\cvEntryClient%
  {Ingeniero Backend Senior}%
  {Globant}%
  {Banco GNB Paraguay}%
  {Plataforma de Banca Digital}%
  {Junio 2022 -- Mayo 2023}

\begin{cvItems}
  \cvItem{Desarrollo de microservicios \textbf{Java / Spring Boot} para gestión de cuentas,
    procesamiento de transacciones y onboarding de clientes de Banco GNB Paraguay.}
  \cvItem{Integración de servicios backend con sistemas bancarios core mediante \textbf{REST APIs}
    y capas de mensajería.}
  \cvItem{Desarrollo orientado a seguridad para cumplimiento regulatorio financiero en el sector
    bancario paraguayo.}
\end{cvItems}

%% ── Experiencia Adicional (condensada) ──────────────────────
\section{Experiencia Adicional}

{\small
\cvMiniEntry{Ingeniero Full Stack}{The Bridge Social (Remoto, EE.UU.)}
\cvMiniEntry{Ingeniero Full Stack}{INDRA\enspace\textbar\enspace Clientes: BCP, RIMAC}
\cvMiniEntry{Ingeniero Full Stack}{Michael Page\enspace\textbar\enspace Cliente: Intercorp}
\cvMiniEntry{Ingeniero Full Stack}{Zoluxiones\enspace\textbar\enspace Cliente: SURA}
\cvMiniEntry{Ingeniero Full Stack}{Experis\enspace\textbar\enspace Cliente: Equifax}
\cvMiniEntry{Ingeniero Full Stack}{Olva Courier}
\cvMiniEntry{Ingeniero Full Stack}{LimaW}
}

%% ── Logros Destacados ────────────────────────────────────────
\section{Logros Destacados}

\begin{cvItems}
  \cvItem{\textbf{Más de 12 años} construyendo sistemas backend con Java, Spring Boot y
    Node.js/NestJS en Salud, Banca, Seguros, Retail y Logística.}
  \cvItem{Arquitectura y entrega de \textbf{AgentCloud}, plataforma de orquestación de agentes IA
    en producción integrando LLMs, Kafka, AWS y pipelines agénticos.}
  \cvItem{Entrega de soluciones backend para clientes de primer nivel: \textbf{Scotiabank},
    SambaSafety, Banco GNB Paraguay, Equifax, Intercorp, SURA, RIMAC y BCP.}
  \cvItem{Liderazgo de equipos de ingeniería distribuidos en proyectos 100\% remotos en
    \textbf{Estados Unidos, América Latina y el Caribe}.}
  \cvItem{Entrega consistente en entornos regulados de alta disponibilidad con estrictos requisitos
    de seguridad y cumplimiento normativo.}
\end{cvItems}
